\section{N-Queens}
\label{sec:rec-n-queens}

\problemstatement{Place $n$ queens on an $n \times n$ chessboard so that
no two queens attack each other (no two share a row, a column, or a
diagonal). Return every distinct valid placement.}

\begin{patternbox}
\emph{``Place $n$ non-attacking items on a board''} is a constraint-satisfaction
backtracking problem where a powerful structural shortcut is available:
since a valid solution can have \textbf{at most one queen per row}
(two queens in the same row always attack each other), we can place
exactly one queen \textbf{per row}, in row order, reducing the search
from ``choose a subset of cells'' to ``choose one column per row,''
checking only column and diagonal conflicts against \emph{already-placed}
queens in earlier rows.
\end{patternbox}

\subsection*{Understanding the problem carefully}

Because we place queens one complete row at a time, by the time we
decide row $r$'s queen, rows $0, \dots, r-1$ already have their queens
correctly and safely placed; we never need to recheck them against each
other again, only check the new candidate column against those already
fixed. A candidate column $c$ in row $r$ is safe if no earlier queen
shares column $c$, and no earlier queen lies on either diagonal passing
through $(r,c)$ -- a queen at $(r', c')$ shares a diagonal with $(r,c)$
exactly when $|r-r'| = |c-c'|$.

\begin{ideabox}[Track Occupied Columns and Diagonals for $O(1)$ Safety Checks]
Maintain three boolean arrays: $\textit{colUsed}[c]$ (is column $c$
occupied by some already-placed queen), and two diagonal arrays indexed
by the constant value along each diagonal direction:
$\textit{diag1}[r+c]$ (the ``$\backslash$''-diagonal) and
$\textit{diag2}[r-c+n-1]$ (the ``$/$''-diagonal, offset to stay
non-negative). $\textsc{Solve}(\textit{row})$: if $\textit{row}=n$:
record the current board; return. Otherwise, for each column $c$ from
$0$ to $n-1$: if none of $\textit{colUsed}[c]$, $\textit{diag1}[r+c]$,
$\textit{diag2}[r-c+n-1]$ are set: place the queen, mark all three
arrays, recurse $\textsc{Solve}(\textit{row}+1)$, then unmark all three
(undo) and remove the queen.
\end{ideabox}

\subsection*{Why diagonal membership reduces to a single sum or difference}

Every cell on the same ``$\backslash$''-diagonal (running from upper-left
to lower-right) shares the same value of $\textit{row}+\textit{col}$;
every cell on the same ``$/$''-diagonal (lower-left to upper-right)
shares the same value of $\textit{row}-\textit{col}$. This means
checking ``does any earlier queen share a diagonal with $(r,c)$'' reduces
to a single array lookup at index $r+c$ (or $r-c$, shifted to stay a
valid non-negative array index) rather than scanning every previously
placed queen and computing $|r-r'|$ versus $|c-c'|$ each time -- turning
an $O(n)$ safety check into an $O(1)$ one, and the whole algorithm's
per-placement cost from $O(n)$ down to $O(1)$.

\subsection*{Algorithm}

\begin{enumerate}
  \item Initialize $\textit{colUsed}, \textit{diag1}, \textit{diag2}$
        as all-false boolean arrays of appropriate sizes.
  \item $\textsc{Solve}(\textit{row})$:
  \begin{enumerate}
    \item If $\textit{row}=n$: record the board; return.
    \item For $c$ from $0$ to $n-1$: if $\textit{colUsed}[c]$,
          $\textit{diag1}[\textit{row}+c]$, and
          $\textit{diag2}[\textit{row}-c+n-1]$ are all false: place the
          queen at $(\textit{row},c)$; set all three arrays true;
          recurse $\textsc{Solve}(\textit{row}+1)$; set all three
          arrays back to false (undo); remove the queen.
  \end{enumerate}
  \item Call $\textsc{Solve}(0)$.
\end{enumerate}

\subsection*{Dry run}

\dryrun{Take $n=4$.}

The search places a queen in row $0$, tries each column, and recurses;
the only two solutions (up to the board's symmetry) for $n=4$ are
columns $[1,3,0,2]$ and $[2,0,3,1]$ (queen in row $0$ at column $1$ or
$2$ respectively, with the rest determined by the safety constraints).
Tracing column $1$ for row $0$: row $1$'s columns $0,1,2$ are all
unsafe (column $1$ taken, or diagonal conflicts), leaving only column
$3$ safe; row $2$'s only safe column, given rows $0,1$, is $0$; row
$3$'s only safe column, given rows $0,1,2$, is $2$ -- yielding the
complete solution $[1,3,0,2]$.

\subsection*{C++ implementation}

\begin{lstlisting}
#include <bits/stdc++.h>
using namespace std;

void solve(int row, int n, vector<int>& cols, vector<bool>& colUsed,
           vector<bool>& diag1, vector<bool>& diag2,
           vector<vector<int>>& result) {
    if (row == n) {
        result.push_back(cols);
        return;
    }
    for (int c = 0; c < n; c++) {
        int d1 = row + c, d2 = row - c + n - 1;
        if (colUsed[c] || diag1[d1] || diag2[d2]) continue;

        cols[row] = c;
        colUsed[c] = diag1[d1] = diag2[d2] = true;

        solve(row + 1, n, cols, colUsed, diag1, diag2, result);

        colUsed[c] = diag1[d1] = diag2[d2] = false;
    }
}

vector<vector<int>> solveNQueens(int n) {
    vector<vector<int>> result;
    vector<int> cols(n, -1);
    vector<bool> colUsed(n, false), diag1(2 * n, false), diag2(2 * n, false);

    solve(0, n, cols, colUsed, diag1, diag2, result);
    return result;
}

int main() {
    auto solutions = solveNQueens(4);
    cout << "Number of solutions: " << solutions.size() << endl;
    for (auto& sol : solutions) {
        for (int c : sol) cout << c << " ";
        cout << endl;
    }
    return 0;
}
\end{lstlisting}

\begin{complexitybox}
\textbf{Time Complexity.} Bounded by $O(n!)$ in the absolute worst case
(one queen per row, up to $n$ column choices each, shrinking as
constraints accumulate), though the $O(1)$ safety check via the
tracking arrays makes each individual placement attempt fast; the true
practical running time is much smaller due to aggressive pruning from
the column/diagonal constraints.
\textbf{Space Complexity.} $O(n)$ for the tracking arrays and the
recursion depth.
\end{complexitybox}

\begin{takeawaybox}
When a backtracking constraint can be reduced to membership in a small
number of equivalence classes (here, columns and the two diagonal
directions), maintain one boolean array per class for $O(1)$
membership checks, instead of recomputing the constraint from scratch
by scanning all previously placed items on every attempt. This
``maintain auxiliary used/unused trackers, mark and unmark them in
lockstep with placing and removing a choice'' idea generalizes directly
to Sudoku's row/column/box constraints in the next section.
\end{takeawaybox}
